\section{Objections and replies}\label{sec:objections}

\subsection{Status functions need collective acceptance, and records do not supply it}
\label{sec:obj-acceptance}

\paragraph{The objection.} Institutional facts exist only as long as they are recognized
\citep{searle1995,searle2010}. A log that nobody accepts creates nothing, so grounding conditions on
records leave out the one thing that makes a social fact social.

\paragraph{Reply.} We do not claim that records replace acceptance. Acceptance attaches to the
constitutive rule, such as the committee's rule that a record signed by three members counts as an
endorsement. Given the rule, each new fact under it is fixed by the records, which is how Searle
treats standing declarations that apply a rule to case after case \citep{searle2010}, and how
\citet{ferraris2013} treats documents that inscribe acts. The sort concerns the facts that a frame,
once accepted, makes depend on records.

\subsection{Grounding is confused with anchoring}\label{sec:obj-anchoring}

\paragraph{The objection.} \citet{epstein2015} separates what makes a social fact hold from what
puts the frame in place. A sort of grounding conditions says nothing about the frame, and the frame
is where the interesting questions of authority live.

\paragraph{Reply.} The sort applies to grounding conditions under fixed frame principles, as
\secref{sec:grounding} states. It also has a corollary for anchoring. Changing a frame asks that no
party still holds the old terms, which is an absence, so a change of frame needs coordination too.
\secref{sec:anchoring} develops what that requires.

\subsection{A signature evidences a fact and does not make it}\label{sec:obj-evidence}

\paragraph{The objection.} Smith presses the point that a record of something is not the thing: an
IOU records a debt and is not itself the debt \citep[p.~289]{smithsearle2003}. A signed record can show that a fact
holds without being what makes it hold.

\paragraph{Reply.} The thesis concerns constitutive records: documents that inscribe an act, in
Ferraris's sense, and what \citet{smith2012} calls document acts. An endorsement, a co-signed seal and
a prescription are acts performed by signing. Many records are evidential, such as the meter's
glucose readings, and the thesis does not make them constitutive of the glucose level. Even an
evidential record grounds institutional facts about itself, such as that this reading was entered by
this meter, and grounds them monotonically.

\subsection{A co-signature is not a joint commitment}\label{sec:obj-joint}

\paragraph{The objection.} Joint commitment, in Gilbert's sense, requires open expressions of
readiness under common knowledge, and creates obligations directed to the other parties. A
cryptographic co-signature may amount to two individual commitments, or to a conditional exchange of
promises. Theorists of shared intention also deny that joint commitment is needed for joint action
\citep{kopec2018}.

\paragraph{Reply.} A co-signed seal meets Gilbert's conditions in the respects that matter here. Each
signature is an open expression of readiness. The record names every co-signer and is shared among
them, which approximates common knowledge. The seal can be reopened only by a further record
co-signed by the same parties, which matches Gilbert's condition that a joint commitment ends only by
the parties together \citep{gilbert2013}. We offer this as a sufficient reading of a co-signed seal as
a joint act, and leave open whether every joint action requires one.

\subsection{Absence facts reduce to positive facts with a timestamp}\label{sec:obj-timestamp}

\paragraph{The objection.} ``No objection by the deadline'' is a positive fact about what a log
contains at a time, so absence facts need no special treatment.

\paragraph{Reply.} Locally this is true, and it is exactly the CALM theorem's point. The fact that a
party has seen no objection by a time is monotone and safe to state.\footnote{In
\texttt{proofs-v1.1.0} a checkpoint that the party signs alone settles what it saw
(\texttt{checkpoint\_positive}), and no set of records that party alone signed holds a seal of a
frame whose closing rule names another party (\texttt{checkpoint\_not\_institutional}).} The institutional fact that no
objection exists is not, because a late or partitioned record can falsify it, and the only way to
know one has heard from everyone is to coordinate \citep{hellerstein2020}. A deadline settles the
institutional fact only if the parties have jointly committed to the deadline as the closing rule,
so the timestamp presupposes the coordination it was meant to remove. Armstrong's totality facts give
the metaphysical parallel: that these are all the facts of a kind is a further fact
\citep[ch.~13]{armstrong1997}.

\subsection{The CALM theorem does not transfer}\label{sec:obj-transfer}

\paragraph{The objection.} The theorem concerns monotone queries over networks of relational
transducers \citep{ameloot2013}. Applying it to social facts is an analogy.

\paragraph{Reply.} The result this paper relies on is proved for its own model, machines over signed
records, and checked in Lean 4 \citep{karna2026,karna2026social}. The theorem has also been generalized beyond
relational transducers, to weaker forms of monotonicity \citep{ameloot2015} and to a semantic
criterion over specifications \citep{hellerstein2026}, and its monotonicity criterion has been applied
to organizational tasks \citep{ju2026}. The claim is stated conditionally throughout: insofar as an institution's facts are
functions of records, the sort applies to them.\footnote{The theorem, \texttt{calm\_social} in
\texttt{proofs-v1.1.0}, carries two conditions. The possible records form a countable set, as
records written as finite strings do. Wherever a fact asks for an absence, the absence shows between
two sets of records that differ by finitely many.}

\subsection{This is the social ontology of blockchains again}\label{sec:obj-blockchain}

\paragraph{The objection.} The social ontology of blockchains already studies how institutional facts
can rest on shared records \citep{reijers2018,davidson2018,suzuki2026}.

\paragraph{Reply.} A blockchain imposes global consensus on every fact. The sort shows that this is
unnecessary for monotone facts, which are most of them. Global consensus is one costly way of
coordinating for absence facts, and a co-signature among the parties a question names is another,
spent only where an absence matters.

\subsection{Institutions are defeasible, and revocation breaks monotonicity}\label{sec:obj-defeasible}

\paragraph{The objection.} Legal and institutional statuses are defeasible \citep{hart1949}: a
license can be revoked and a decision appealed \citep{prakken2004}. A theory built on monotone records
cannot represent them.

\paragraph{Reply.} The log stays monotone, and derived statuses such as a valid license or a current
dose may not be. Revocation is precisely an absence-dependent status, since a license is valid only
if no revocation of it exists. Defeasibility therefore confirms the sort, and a revocation list that
is published, signed and dated is the seal that lets a verifier say valid as of that list.\footnote{\texttt{valid\_asks\_absence}
and \texttt{valid\_asOf\_positive} in \texttt{proofs-v1.1.0}, and for a rule kept as records that a
later record can repeal, \texttt{inForce\_asks\_absence}.}
